% =====================================================================
%  style.tex -- the visual language of the book (see STYLE.md)
%  Loaded after common.tex.  Every chapter uses only these environments:
%    theorem-like (boxed):  theorem, proposition, lemma, corollary, definition
%    worked examples:       example
%    pedagogical boxes:     intuition, keyidea, pitfall
%    chapter frame:         chapterintro (+ \prerequisites), chaptersummary
%    exercises:             exercise (in an unnumbered "Exercises" section)
% =====================================================================

% ---------- fonts ------------------------------------------------------
\usepackage[sc]{mathpazo}
\linespread{1.06}

% ---------- palette ----------------------------------------------------
\definecolor{bookaccent}{RGB}{31,78,121}     % headings, theorem frames
\definecolor{bookdef}{RGB}{34,110,72}        % definitions
\definecolor{bookexample}{RGB}{176,96,16}    % examples
\definecolor{bookintuition}{RGB}{70,110,170} % intuition boxes
\definecolor{bookkey}{RGB}{120,60,140}       % key ideas
\definecolor{bookpitfall}{RGB}{170,40,40}    % pitfalls
\definecolor{booksummary}{RGB}{60,60,60}     % chapter summary

% ---------- headings -----------------------------------------------------
\usepackage{titlesec}
\titleformat{\chapter}[display]
  {\normalfont\color{bookaccent}}
  {\scshape\large Chapter \thechapter}
  {6pt}
  {\Huge\bfseries\color{black}}
  [\vspace{2pt}{\color{bookaccent}\titlerule[1pt]}]
\titlespacing*{\chapter}{0pt}{0pt}{18pt}
\titleformat{\section}
  {\normalfont\Large\bfseries}{\color{bookaccent}\thesection}{0.8em}{}
\titleformat{\subsection}
  {\normalfont\large\bfseries}{\color{bookaccent}\thesubsection}{0.8em}{}
\titleformat{\subsubsection}
  {\normalfont\normalsize\bfseries}{\thesubsubsection}{0.8em}{}
\titleformat{\paragraph}[runin]
  {\normalfont\normalsize\bfseries}{\theparagraph}{0.5em}{}

% subtitle under the chapter title
\renewcommand{\chaptersubtitle}[1]{%
  \vspace*{-10pt}{\large\itshape\color{bookaccent!80!black} #1\par}\vspace{14pt}}

% ---------- boxed theorem-like environments ------------------------------
\tcbset{
  bookthm/.style={enhanced, breakable, boxrule=0pt, frame hidden, sharp corners,
    borderline west={2.5pt}{0pt}{#1}, colback=#1!5!white,
    left=8pt, right=6pt, top=4pt, bottom=4pt, before skip=10pt, after skip=10pt},
  bookbox/.style={enhanced, breakable, sharp corners, boxrule=0.6pt,
    colframe=#1, colback=#1!4!white, coltitle=white, colbacktitle=#1,
    fonttitle=\bfseries\small, left=8pt, right=8pt, top=5pt, bottom=5pt,
    before skip=10pt, after skip=10pt, attach boxed title to top left={xshift=8pt, yshift=-2mm},
    boxed title style={sharp corners, boxrule=0pt}},
}
\tcolorboxenvironment{theorem}{bookthm=bookaccent}
\tcolorboxenvironment{proposition}{bookthm=bookaccent}
\tcolorboxenvironment{prop}{bookthm=bookaccent}
\tcolorboxenvironment{lemma}{bookthm=bookaccent}
\tcolorboxenvironment{corollary}{bookthm=bookaccent}
\tcolorboxenvironment{claim}{bookthm=bookaccent}
\tcolorboxenvironment{definition}{bookthm=bookdef}
\tcolorboxenvironment{assumption}{bookthm=bookdef}
\tcolorboxenvironment{example}{bookthm=bookexample}

% ---------- book-wide notation (see STYLE.md) ------------------------------
\DeclareMathOperator{\Var}{Var}
\DeclareMathOperator{\Cov}{Cov}
\newcommand{\Xs}{\mathcal{X}}          % input space
\newcommand{\Ys}{\mathcal{Y}}          % output space
\newcommand{\Hs}{\mathcal{F}}          % hypothesis class
\newcommand{\Risk}{R}                  % population risk R(f)
\newcommand{\ERisk}{\widehat{R}_n}     % empirical risk on n points
\newcommand{\given}{\,\middle|\,}      % conditioning bar inside \left( \right)

% ---------- pedagogical boxes ---------------------------------------------
\newtcolorbox{intuition}[1][Intuition]{bookbox=bookintuition, title={#1}}
\newtcolorbox{keyidea}[1][Key idea]{bookbox=bookkey, title={#1}}
\newtcolorbox{pitfall}[1][Pitfall]{bookbox=bookpitfall, title={#1}}

% ---------- chapter frame --------------------------------------------------
\newcommand{\prerequisites}[1]{%
  \par\smallskip\noindent{\color{bookaccent}\textbf{Prerequisites.}} #1\par}
\newtcolorbox{chapterintro}{enhanced, breakable, sharp corners, boxrule=0pt,
  frame hidden, colback=bookaccent!6!white, borderline north={1pt}{0pt}{bookaccent},
  borderline south={1pt}{0pt}{bookaccent}, left=8pt, right=8pt, top=6pt, bottom=6pt,
  before upper={\noindent{\color{bookaccent}\textbf{In this chapter}}\par\smallskip}}
\newtcolorbox{chaptersummary}{bookbox=booksummary, title={Summary}}

% exercises are numbered within the chapter, independently of theorems
\newcommand{\hint}[1]{\par\noindent{\small\itshape Hint: #1}}

% ---------- solutions to exercises (collected in backmatter/solutions.tex) ----
\newcommand{\solutionlink}[1]{\par\nopagebreak\hfill{\small\hyperref[sol:#1]{Solution, p.~\pageref*{sol:#1}}}}
\newcommand{\solution}[1]{\par\medskip\noindent\textbf{Exercise~\ref*{#1}}\label{sol:#1}%
  \hfill{\small\hyperref[#1]{statement, p.~\pageref*{#1}}}\par\nopagebreak\smallskip}
